\chapter*{Післямова}
\addcontentsline{toc}{chapter}{Післямова}

Вихід цього посібника став результатом плідної праці, глибоких досліджень та професійної підтримки багатьох фахівців, чий досвід дозволив перетворити теоретичні знання на дієвий інструмент захисту.

Насамперед щиро дякую моїм рідним та близьким — за ваші віру, терпіння, повсякденну турботу й невичерпну мотивацію, які надихали мене та допомагали долати всі труднощі на цьому шляху.

Особливу вдячність висловлюю особливу вдячність \textbf{Українській міжбанківській Асоціації членів платіжних систем «ЄМА»}. Ваші рекомендації, аналітичні дані та верифікація матеріалу стали надійним фундаментом для цієї книги, гарантуючи відповідність викладених алгоритмів передовим стандартам протидії шахрайству.

Робота набула необхідної глибини завдяки моїм рецензентам та консультантам. Щиро дякую \textbf{І.C. ~Петровій}, \textbf{Р.~Федоровській} та \textbf{С.О.~Гнатюку} за експертну оцінку технічних та безпекових аспектів, а також \textbf{І.А.~Сліпухіній} за методичні поради. Окрему подяку висловлюю \textbf{С.А.~Данилову} та \textbf{І.І.~Чубатюк}: ваші фахові роз’яснення у сферах нейробіології та психології дозволили розкрити «людський фактор» кібершахрайства і пояснити, чому наш мозок потрапляє у пастки зловмисників.

Важливу роль у технічній реалізації досліджень відіграв \textbf{П.В.~Таранущенко}. Його допомога у проведенні аудитів безпеки в законних цілях та моделюванні технічних сценаріїв атак дозволила перевірити на практиці теоретичні гіпотези, які неможливо протестувати самотужки.

Дякую колегам і друзям за підтримку, пошук практичних ситуацій та дискусії, що ставали джерелом натхнення.

\vspace{0.5cm}
\centerline{\rule{0.5\linewidth}{0.5pt}} % Розділова лінія
\vspace{0.1cm}

\clearpage

\chapter *{Контакти автора}

\begin{center}
    
    З питань співпраці, консультацій та доступу до CyberSimulator, будь ласка, звертайтеся за наступними контактами:
    
    \vspace{0.1cm}
    \label{author:contacts:preface}
    
    \textbf{E-mail (основний):} \href{mailto:s.yevdokymov@ukr.net}{s.yevdokymov@ukr.net} \\
    \vspace{0.1cm}
    \textbf{E-mail (додатковий):} \href{mailto:serhii.yevdokymov2022@gmail.com}{serhii.yevdokymov2022@gmail.com} \\
    \vspace{0.1cm}
   %  \textbf{Telegram:} \href{https://t.me/sargeras197}{@sargeras197} \\
   %  \vspace{0.1cm}

  %  \vspace{0.1cm}
  %  \textbf{ORCID:} \href{https://orcid.org/0000-0001-7213-0259}{0000-0001-7213-0259}

\end{center}

   % \vspace{0.2cm}

QR-код для швидкого доступу до контактів автора:

\begin{figure}[h]
    \centering
    % Переконайтеся, що файл fig1_1.jpg (або png) є у папці img
    \includegraphics[width=0.7\linewidth]{img/qr-code-contacts.png} 
    %\caption{Збитки від шахрайств з платіжними картками в Україні}
    \label{fig:1_qr_contacts}
\end{figure}

 %\vspace{0.2cm}
 
Платформа надає можливiсть ознайомлення з теоретичними матерiалами з протидiї кiбершахрайству, мiстить практичнi завдання, iнтерактивнi сценарiї та навчальнi челенджi, що моделюють типовi кiберзагрози в цифровому середовищi. Вiдкритий формат доступу дозволяє широкому колу користувачiв брати участь у навчальному процесi, застосовувати отриманi знання на практицi та сформувати навички безпечної поведiнки в цифровому просторi.

\clearpage % Ця команда створює розрив сторінки, щоб контакти були на новій сторінці
